% TOP_PROBLEM: {"id": 2622, "title": "Kourovka Notebook Problem 21.113", "category_id": 17, "category": {"id": 17, "name": "group_theory", "display_name": "Group Theory", "description": "Problems about groups, group actions, representations, and related algebraic structures.", "slug": "group-theory", "order_index": 17, "created_at": "2026-07-31T15:26:25.671Z"}, "status": "open", "problem_number": "KOU-21.113", "set_id": 10}
% FIRST_POSTED: 2026-10-01
% ENTRY_KIND: statement_only
% UnsolvedMath ID 2622; source catalogue code KOU-21.113
% !TeX program = lualatex
% Definition and sources only; this file is not an LLM solution attempt.
\documentclass[11pt,a4paper]{article}
\usepackage[margin=25mm]{geometry}
\usepackage{fontspec}
\setmainfont{lmroman10-regular.otf}[BoldFont=lmroman10-bold.otf,
 ItalicFont=lmroman10-italic.otf,BoldItalicFont=lmroman10-bolditalic.otf]
\usepackage{amsmath,amssymb,mathtools}
\usepackage{unicode-math}
\setmathfont{latinmodern-math.otf}
\AtBeginDocument{\let\setminus\smallsetminus}
\usepackage{xurl,hyperref}
\hypersetup{colorlinks=true,urlcolor=blue}
\setcounter{secnumdepth}{0}
\setlength{\parindent}{0pt}
\setlength{\parskip}{5pt}
\setlength{\emergencystretch}{3em}
\begin{document}
\section*{Kourovka Notebook Problem 21.113}\label{2622}
{\small UnsolvedMath ID 2622 · KOU-21.113 · Group Theory · Kourovka Notebook - New Problems, Issue 21}
\subsection{Definitions and mathematical statement}
Let $G$ be finite and $p$ prime. An element is $p$-regular if its order is prime to $p$, and $p$-singular otherwise. A $p$-element has order a power of $p$, including the identity of order $p^0$. Set $C_G(x)=\{g:gx=xg\}$ and define the conjugacy-invariant function
\[
 \Psi_{p,G}(x)=
 \begin{cases}
 |\{y\in C_G(x):|y|\text{ is a power of }p\}|,
       &x\text{ is }p\text{-regular},\\
 0,&x\text{ is }p\text{-singular}.
 \end{cases}
\]
\textbf{(a)} Is $\Psi_{p,G}$ always an ordinary character of $G$? This means the trace character of a finite-dimensional characteristic-zero representation, or equivalently a sum of irreducible complex characters with nonnegative integer coefficients.

\textbf{(b)} If so, can it be afforded by a finitely generated projective $RG$-module? Here $R$ is a complete discrete valuation ring of characteristic zero, with maximal ideal $J(R)$, fraction field $K$, and residue field $k=R/J(R)$ of characteristic $p$. Require $K$ and $k$ to be splitting fields for $G$ and all its subgroups: simple representations over either field remain simple after extending scalars to an algebraic closure.

A projective $RG$-module is a direct summand of a finite-rank free $RG$-module. To afford the function means that the $KG$-module $K\otimes_R P$ has trace $\Psi_{p,G}(x)$ at every $x$. The questions are universal in $G$ and $p$, with the indicated splitting-ring hypotheses in the stronger integral realization.
\subsection{Short English statement}
Is this centralizer-counting function always an ordinary character, and can it moreover be realized by a projective module over a splitting valuation ring?
\subsection{Sources}
\begin{itemize}
\item[S1] ULAM.AI and the UnsolvedMath Contributors, supplied problems.json, record 2622 (KOU-21.113), Kourovka Notebook - New Problems, Issue 21. Catalogue identity and source transcription; CC BY 4.0.
\url{https://huggingface.co/datasets/ulamai/UnsolvedMath}.
\item[S2] The Kourovka Notebook, 21st edition, version 47 (30 September 2026), new problems, Problem 21.113. Expanded formulation of the supplied catalogue statement.
\url{https://arxiv.org/pdf/1401.0300v47}.
\end{itemize}
\end{document}
